\chapter{L9 -- Analisi modale e stabilità a partire dalla forma di Jordan}

\section{Obiettivo della lezione}

Dopo aver introdotto la forma di Jordan e le catene di autovettori generalizzati, il passo naturale successivo è capire \emph{che cosa si vede nel tempo}.

L'idea fondamentale è questa:

\begin{quote}
se porto la matrice di stato in forma di Jordan, allora nell'evoluzione libera del sistema compaiono in modo esplicito tutti i possibili andamenti temporali.
\end{quote}

Questi andamenti temporali elementari vengono chiamati \textbf{modi del sistema}.  
La struttura dei modi dipende da:
\begin{itemize}
    \item il valore degli autovalori;
    \item il fatto che siano reali oppure complessi;
    \item la dimensione del blocco di Jordan più grande associato a ciascun autovalore.
\end{itemize}

Questa lezione serve quindi a collegare in modo chiaro:
\[
\text{autovalori e blocchi di Jordan}
\qquad \Longleftrightarrow \qquad
\text{modi temporali del sistema}
\qquad \Longleftrightarrow \qquad
\text{stabilità / instabilità}.
\]

\section{Perché i blocchi di Jordan sono importanti}

Nel tempo continuo, l'evoluzione libera è
\[
x(t)=e^{At}x(0),
\]
mentre nel tempo discreto è
\[
x(k)=A^k x(0).
\]

Se $A$ è simile a una forma di Jordan $J$, cioè
\[
A=QJQ^{-1},
\]
allora
\[
e^{At}=Qe^{Jt}Q^{-1}
\qquad \text{(tempo continuo)},
\]
e
\[
A^k=QJ^kQ^{-1}
\qquad \text{(tempo discreto)}.
\]

Questo significa che:
\begin{itemize}
    \item le matrici $Q$ e $Q^{-1}$ fanno solo un cambio di coordinate;
    \item i veri andamenti nel tempo compaiono dentro $e^{Jt}$ oppure dentro $J^k$.
\end{itemize}

Di conseguenza, per fare analisi modale, bisogna guardare i blocchi di Jordan.

\section{Modi critici: perché compaiono i termini polinomiali}

Supponiamo di avere, nel continuo, un blocco di Jordan reale associato a un autovalore critico.  
Se ad esempio compare il blocco
\[
J_0^{(3)}=
\begin{pmatrix}
0 & 1 & 0\\
0 & 0 & 1\\
0 & 0 & 0
\end{pmatrix},
\]
allora
\[
e^{J_0^{(3)}t}
=
I + J_0^{(3)}t + \frac{(J_0^{(3)})^2 t^2}{2!}
=
\begin{pmatrix}
1 & t & \dfrac{t^2}{2}\\
0 & 1 & t\\
0 & 0 & 1
\end{pmatrix}.
\]

Quindi i modi che compaiono sono:
\[
1,\qquad t,\qquad t^2.
\]

\paragraph{Interpretazione.}
\begin{itemize}
    \item Il modo $1$ è un andamento costante: non esplode, ma nemmeno converge a zero.
    \item Il modo $t$ diverge lentamente, in modo polinomiale.
    \item Il modo $t^2$ diverge ancora più rapidamente.
\end{itemize}

Per questo motivo, quando un autovalore è \emph{critico} (nel continuo $\Re(\lambda)=0$, nel discreto $|\lambda|=1$), non basta conoscere il suo valore: bisogna sapere anche la dimensione del blocco di Jordan più grande a esso associato.

\begin{figure}[H]
\centering
\begin{tikzpicture}[>=Latex, node distance=1.4cm]
    \node[draw, rounded corners, align=center] (b1) {$J^{(1)}_0$};
    \node[draw, rounded corners, align=center, right=of b1] (m1) {$1$};

    \node[draw, rounded corners, align=center, below=1.1cm of b1] (b2) {$J^{(2)}_0$};
    \node[draw, rounded corners, align=center, right=of b2] (m2) {$1,\; t$};

    \node[draw, rounded corners, align=center, below=1.1cm of b2] (b3) {$J^{(3)}_0$};
    \node[draw, rounded corners, align=center, right=of b3] (m3) {$1,\; t,\; t^2$};

    \draw[->, thick] (b1) -- (m1);
    \draw[->, thick] (b2) -- (m2);
    \draw[->, thick] (b3) -- (m3);
\end{tikzpicture}
\caption{Un blocco di Jordan critico di ordine crescente genera modi polinomiali di grado crescente.}
\end{figure}

\section{Esempio strutturale: lettura dei modi da una matrice a blocchi}

Consideriamo una matrice a blocchi triangolare superiore del tipo
\begin{equation}
A=
\begin{pmatrix}
A_1^\star & B_{12}\\
0 & A_2
\end{pmatrix},
\label{eq:L9_Ablock}
\end{equation}
con
\[
A_1^\star=
\begin{pmatrix}
1 & 0 & 0\\
-1 & -1 & 1\\
-3 & -1 & -1
\end{pmatrix},
\qquad
A_2=
\begin{pmatrix}
1 & 0 & 0\\
-1 & -2 & 0\\
0 & 1 & -2
\end{pmatrix}.
\]

Il blocco di accoppiamento $B_{12}$ può essere non nullo, ma gli autovalori della matrice triangolare a blocchi sono comunque l'unione degli autovalori dei blocchi diagonali.

\subsection{Autovalori dei blocchi diagonali}

Per $A_1^\star$ si riconoscono:
\begin{itemize}
    \item un autovalore reale $\lambda=1$;
    \item una coppia complessa coniugata $\lambda=-1\pm j$.
\end{itemize}

Per $A_2$ si hanno:
\begin{itemize}
    \item un autovalore reale $\lambda=1$;
    \item l'autovalore $\lambda=-2$ con molteplicità algebrica $2$.
\end{itemize}

Quindi, per l'intera matrice $A$, gli autovalori sono:
\[
\lambda_1 = 1 \quad \text{con } m_a(1)=2,
\]
\[
\lambda_2 = -2 \quad \text{con } m_a(-2)=2,
\]
\[
\lambda_{3,4} = -1\pm j \quad \text{con molteplicità }1.
\]

\paragraph{Primo commento importante.}
Sapere gli autovalori non basta ancora.  
Dobbiamo capire \emph{quanto sono grandi i blocchi di Jordan} associati a ciascuno di essi.

\subsection{Autovalore \texorpdfstring{$\lambda=1$}{lambda=1}}

Dallo studio del rango di $A-I$ si trova
\[
\rank(A-I)=5,
\]
quindi
\[
\mu_1(1)=\dim\ker(A-I)=6-5=1.
\]

Poiché la molteplicità algebrica di $\lambda=1$ è $2$, ma la molteplicità geometrica è $1$, segue che per $\lambda=1$ c'è \emph{un solo blocco di Jordan}, necessariamente di ordine $2$:
\[
J^{(2)}_{1}=
\begin{pmatrix}
1 & 1\\
0 & 1
\end{pmatrix}.
\]

\subsection{Autovalore \texorpdfstring{$\lambda=-2$}{lambda=-2}}

In modo analogo, studiando $A+2I$ si trova
\[
\rank(A+2I)=5,
\]
quindi
\[
\mu_1(-2)=\dim\ker(A+2I)=1.
\]

Anche qui la molteplicità algebrica è $2$ ma la molteplicità geometrica è $1$, dunque per $\lambda=-2$ compare un solo blocco di Jordan di ordine $2$:
\[
J^{(2)}_{-2}=
\begin{pmatrix}
-2 & 1\\
0 & -2
\end{pmatrix}.
\]

\subsection{Coppia complessa \texorpdfstring{$-1\pm j$}{-1 +- j}}

La coppia complessa coniugata $-1\pm j$ compare una sola volta, quindi nel formalismo reale è associata a un unico blocco reale $2\times 2$:
\[
M=
\begin{pmatrix}
-1 & 1\\
-1 & -1
\end{pmatrix}.
\]

\paragraph{Conclusione.}
La forma di Jordan reale della matrice avrà la struttura
\begin{equation}
J=
\operatorname{diag}
\left(
J_1^{(2)},
J_{-2}^{(2)},
M
\right),
\label{eq:L9_J_example}
\end{equation}
a meno di una diversa permutazione dei blocchi sulla diagonale.

\section{Modi che compaiono nell'esempio}

Una volta trovata la struttura di Jordan \eqref{eq:L9_J_example}, i modi si leggono immediatamente.

\subsection{Contributo del blocco \texorpdfstring{$J_1^{(2)}$}{J1}}

Nel continuo:
\[
e^{J_1^{(2)}t}
=
e^t
\begin{pmatrix}
1 & t\\
0 & 1
\end{pmatrix},
\]
quindi i modi associati sono
\[
e^t,
\qquad
t e^t.
\]

\subsection{Contributo del blocco \texorpdfstring{$J_{-2}^{(2)}$}{J-2}}

Si ha
\[
e^{J_{-2}^{(2)}t}
=
e^{-2t}
\begin{pmatrix}
1 & t\\
0 & 1
\end{pmatrix},
\]
quindi i modi associati sono
\[
e^{-2t},
\qquad
t e^{-2t}.
\]

\subsection{Contributo del blocco complesso reale \texorpdfstring{$M$}{M}}

Per il blocco
\[
M=
\begin{pmatrix}
-1 & 1\\
-1 & -1
\end{pmatrix}
=
- I + 
\begin{pmatrix}
0 & 1\\
-1 & 0
\end{pmatrix},
\]
si ottiene
\[
e^{Mt}
=
e^{-t}
\begin{pmatrix}
\cos t & \sin t\\
-\sin t & \cos t
\end{pmatrix}.
\]

Dunque i modi associati alla coppia complessa sono:
\[
e^{-t}\cos t,
\qquad
e^{-t}\sin t.
\]

\paragraph{Riassumendo.}
I modi dell'esempio sono:
\begin{equation}
e^t,\qquad te^t,\qquad e^{-2t},\qquad te^{-2t},\qquad e^{-t}\cos t,\qquad e^{-t}\sin t.
\label{eq:L9_modes_example}
\end{equation}

\paragraph{Interpretazione.}
\begin{itemize}
    \item $e^t$ e $te^t$ divergono;
    \item $e^{-2t}$ e $te^{-2t}$ convergono a zero;
    \item $e^{-t}\cos t$ ed $e^{-t}\sin t$ sono oscillazioni smorzate.
\end{itemize}

Quindi il sistema dell'esempio è instabile, perché compaiono modi divergenti.

\section{Osservazione importante: non sempre serve costruire tutte le catene}

Se l'obiettivo è capire gli \emph{andamenti nel tempo}, spesso non è necessario costruire esplicitamente tutte le catene di autovettori generalizzati.

Infatti:
\begin{itemize}
    \item per sapere quali modi compaiono, basta conoscere gli autovalori e la dimensione del blocco di Jordan più grande a essi associato;
    \item la costruzione delle catene serve invece quando si vuole anche esplicitare il cambio di coordinate $Q$ e quindi scrivere in dettaglio le componenti dello stato nelle coordinate originali.
\end{itemize}

\section{Esempio: come ottenere una traiettoria costante}

Consideriamo il sistema continuo
\begin{equation}
\dot{x}(t)=Ax(t),
\qquad
A=
\begin{pmatrix}
6 & -3\\
2 & -1
\end{pmatrix}.
\label{eq:L9_Asmall}
\end{equation}

Vogliamo trovare $x(0)$ tale che la traiettoria sia costante, cioè
\[
x(t)=
\begin{pmatrix}
k_1\\
k_2
\end{pmatrix}
\qquad \forall t.
\]

\subsection*{Autovalori}

Si ha
\[
\det(A-\lambda I)=
\det\begin{pmatrix}
6-\lambda & -3\\
2 & -1-\lambda
\end{pmatrix}
=
(6-\lambda)(-1-\lambda)+6
=
\lambda(\lambda-5).
\]

Quindi gli autovalori sono
\[
\lambda_1=0,
\qquad
\lambda_2=5.
\]

\subsection*{Autovettori}

Per $\lambda=0$:
\[
Ax=0
\quad \Longrightarrow \quad
\ker A = \operatorname{span}\left\{
\begin{pmatrix}
1\\
2
\end{pmatrix}
\right\}.
\]

Per $\lambda=5$:
\[
(A-5I)x=0
\quad \Longrightarrow \quad
\ker(A-5I)=\operatorname{span}\left\{
\begin{pmatrix}
3\\
1
\end{pmatrix}
\right\}.
\]

Possiamo quindi scegliere
\[
Q=
\begin{pmatrix}
1 & 3\\
2 & 1
\end{pmatrix},
\qquad
J=
\begin{pmatrix}
0 & 0\\
0 & 5
\end{pmatrix}.
\]

Allora
\[
e^{At}=Qe^{Jt}Q^{-1}
=
Q
\begin{pmatrix}
1 & 0\\
0 & e^{5t}
\end{pmatrix}
Q^{-1}.
\]

\paragraph{Idea chiave.}
Per ottenere una traiettoria costante, devo eliminare il modo $e^{5t}$ e lasciare solo il modo associato all'autovalore nullo.

Questo accade se la condizione iniziale appartiene all'autospazio di $\lambda=0$, cioè se
\[
x(0)=\alpha
\begin{pmatrix}
1\\
2
\end{pmatrix}.
\]

\paragraph{Conclusione.}
Tutte e sole le condizioni iniziali del tipo
\begin{equation}
x(0)=\alpha
\begin{pmatrix}
1\\
2
\end{pmatrix},
\qquad \alpha\in\R,
\label{eq:L9_constant_traj}
\end{equation}
generano traiettorie costanti.

\paragraph{Interpretazione modale.}
Se scelgo una condizione iniziale nel sottospazio associato al modo $1=e^{0t}$, allora lo stato non vede il modo esponenziale $e^{5t}$ e rimane costante nel tempo.

\section{Analisi modale generale nel tempo continuo}

\subsection{Autovalore reale}

Sia $\lambda_i \in \R$ un autovalore e sia $p_i$ la dimensione del blocco di Jordan \emph{più grande} associato a $\lambda_i$.

Allora i modi associati a $\lambda_i$ sono del tipo
\begin{equation}
e^{\lambda_i t},\qquad
t e^{\lambda_i t},\qquad
t^2 e^{\lambda_i t},\qquad
\dots,\qquad
t^{p_i-1} e^{\lambda_i t}.
\label{eq:L9_modes_real_CT}
\end{equation}

\paragraph{Commento.}
\begin{itemize}
    \item Se $p_i=1$, compare solo il modo esponenziale puro.
    \item Se $p_i>1$, compaiono anche fattori polinomiali in $t$.
\end{itemize}

\subsection{Autovalore complesso}

Se invece
\[
\lambda_i = \sigma_i + j\omega_i,
\qquad
\overline{\lambda_i} = \sigma_i - j\omega_i,
\]
allora nel formalismo reale i modi compaiono a coppie:
\begin{equation}
e^{\sigma_i t}\cos(\omega_i t),
\qquad
e^{\sigma_i t}\sin(\omega_i t).
\label{eq:L9_modes_complex_CT_base}
\end{equation}

Se il blocco reale di Jordan associato ha dimensione $q_i$ (equivalentemente, se nel complesso il blocco associato ha ordine $q_i$), allora compaiono anche i fattori polinomiali:
\begin{equation}
e^{\sigma_i t}\cos(\omega_i t),\;
e^{\sigma_i t}\sin(\omega_i t),\;
t e^{\sigma_i t}\cos(\omega_i t),\;
t e^{\sigma_i t}\sin(\omega_i t),\;
\dots,\;
t^{q_i-1} e^{\sigma_i t}\cos(\omega_i t),\;
t^{q_i-1} e^{\sigma_i t}\sin(\omega_i t).
\label{eq:L9_modes_complex_CT}
\end{equation}

\paragraph{Interpretazione.}
Questi modi rappresentano oscillazioni:
\begin{itemize}
    \item smorzate se $\sigma_i<0$;
    \item amplificate se $\sigma_i>0$;
    \item non smorzate se $\sigma_i=0$.
\end{itemize}

\begin{figure}[H]
\centering
\begin{tikzpicture}[>=Latex, node distance=1.25cm]
    \node[draw, rounded corners, align=center] (r1) {$J_{\lambda}^{(1)}$\\$\lambda\in\R$};
    \node[draw, rounded corners, align=center, right=3.8cm of r1] (m1) {$e^{\lambda t}$};

    \node[draw, rounded corners, align=center, below=1.2cm of r1] (r2) {$J_{\lambda}^{(3)}$\\$\lambda\in\R$};
    \node[draw, rounded corners, align=center, right=3.8cm of r2] (m2) {$e^{\lambda t},\; t e^{\lambda t},\; t^2 e^{\lambda t}$};

    \node[draw, rounded corners, align=center, below=1.2cm of r2] (c1) {$\text{blocco reale }M$\\$\lambda=\sigma\pm j\omega$};
    \node[draw, rounded corners, align=center, right=3.8cm of c1] (m3) {$e^{\sigma t}\cos(\omega t),\; e^{\sigma t}\sin(\omega t)$};

    \draw[->, thick] (r1) -- (m1);
    \draw[->, thick] (r2) -- (m2);
    \draw[->, thick] (c1) -- (m3);
\end{tikzpicture}
\caption{Ogni blocco di Jordan genera una famiglia caratteristica di modi temporali.}
\end{figure}

\section{Analisi modale generale nel tempo discreto}

Nel discreto, se
\[
x(k+1)=Ax(k),
\]
allora
\[
x(k)=A^k x(0)=QJ^kQ^{-1}x(0).
\]

\subsection{Autovalore reale}

Se $\lambda_i\in\R$ e il blocco di Jordan più grande associato ha dimensione $p_i$, allora i modi sono del tipo
\begin{equation}
\lambda_i^k,\qquad
k\lambda_i^k,\qquad
k^2\lambda_i^k,\qquad
\dots,\qquad
k^{p_i-1}\lambda_i^k,
\label{eq:L9_modes_real_DT}
\end{equation}
a meno di costanti moltiplicative o equivalenze del tipo $\binom{k}{r}\lambda_i^{k-r}$.

\paragraph{Commento.}
Per l'analisi qualitativa, quello che conta davvero è la presenza del fattore polinomiale in $k$.

\subsection{Autovalore complesso}

Se
\[
\lambda_i=\sigma_i+j\omega_i=\rho_i e^{j\theta_i},
\qquad
\overline{\lambda_i}=\rho_i e^{-j\theta_i},
\]
allora i modi reali di base sono
\begin{equation}
\rho_i^k \cos(k\theta_i),
\qquad
\rho_i^k \sin(k\theta_i).
\label{eq:L9_modes_complex_DT_base}
\end{equation}

Se il blocco reale di Jordan più grande associato alla coppia ha ordine $q_i$, allora compaiono anche i fattori polinomiali:
\begin{equation}
\rho_i^k \cos(k\theta_i),\;
\rho_i^k \sin(k\theta_i),\;
k\rho_i^k \cos(k\theta_i),\;
k\rho_i^k \sin(k\theta_i),\;
\dots,\;
k^{q_i-1}\rho_i^k \cos(k\theta_i),\;
k^{q_i-1}\rho_i^k \sin(k\theta_i).
\label{eq:L9_modes_complex_DT}
\end{equation}

\paragraph{Osservazione.}
In formule più precise, nei blocchi discreti possono comparire anche potenze del tipo $\rho_i^{k-r}$; tuttavia, dal punto di vista asintotico, la struttura è quella di un termine oscillatorio-amplificato/smorzato moltiplicato per un polinomio in $k$.

\section{Tutti questi andamenti sono i modi del sistema}

Possiamo ora formulare chiaramente il concetto.

\paragraph{Definizione operativa.}
Le funzioni di tempo che compaiono nelle componenti di
\[
e^{Jt}
\qquad \text{oppure} \qquad
J^k
\]
sono i \textbf{modi propri del sistema}.

Nel continuo i modi sono funzioni di $t$; nel discreto sono successioni in $k$.

\paragraph{Messaggio chiave.}
Una volta nota la forma di Jordan, i modi si leggono direttamente e quindi si può prevedere qualitativamente l'evoluzione dello stato senza risolvere ogni volta il sistema da zero.

\section{Stabilità letta direttamente dai modi}

Adesso usiamo i modi per capire quando il sistema converge, diverge oppure resta limitato.

\subsection{Tempo continuo}

Sia $\lambda_i$ un autovalore di $A$.

\paragraph{Caso 1: \texorpdfstring{$\Re(\lambda_i)<0$}{Re(lambda)<0}.}
Tutti i modi associati hanno il fattore
\[
e^{\Re(\lambda_i)t},
\]
che tende a zero esponenzialmente. Anche se compaiono fattori polinomiali $t^r$, il decadimento esponenziale domina, quindi i modi convergono a zero.

\paragraph{Caso 2: \texorpdfstring{$\Re(\lambda_i)>0$}{Re(lambda)>0}.}
I modi divergono esponenzialmente. Eventuali fattori polinomiali rendono solo più marcata la divergenza, ma il comportamento qualitativo è già deciso dal segno positivo della parte reale.

\paragraph{Caso 3: \texorpdfstring{$\Re(\lambda_i)=0$}{Re(lambda)=0}.}
Questo è il caso critico.

\begin{itemize}
    \item Se il blocco di Jordan massimo associato a $\lambda_i$ ha dimensione $1$, allora compaiono solo modi costanti oppure oscillatori limitati, per esempio
    \[
    1,\qquad \cos(\omega t),\qquad \sin(\omega t).
    \]
    \item Se invece il blocco massimo ha dimensione maggiore di $1$, compaiono modi con fattori polinomiali:
    \[
    t,\; t^2,\; \dots,\; t^r,
    \qquad
    t\cos(\omega t),\; t\sin(\omega t),\; \dots
    \]
    e quindi il sistema diventa polinomialmente divergente.
\end{itemize}

\paragraph{Riassunto per il continuo.}
\begin{itemize}
    \item $\Re(\lambda_i)<0$ $\Rightarrow$ modi convergenti;
    \item $\Re(\lambda_i)>0$ $\Rightarrow$ modi divergenti;
    \item $\Re(\lambda_i)=0$ $\Rightarrow$ bisogna guardare la dimensione del blocco di Jordan massimo.
\end{itemize}

\subsection{Tempo discreto}

Nel discreto il ruolo di $\Re(\lambda_i)$ è svolto dal modulo $|\lambda_i|$.

\paragraph{Caso 1: \texorpdfstring{$|\lambda_i|<1$}{|lambda|<1}.}
Tutti i modi convergono a zero. Se $\lambda_i$ è complesso, possono esserci oscillazioni, ma con ampiezza decrescente.

\paragraph{Caso 2: \texorpdfstring{$|\lambda_i|>1$}{|lambda|>1}.}
I modi divergono esponenzialmente. Se $\lambda_i$ è complesso, la divergenza può essere oscillatoria.

\paragraph{Caso 3: \texorpdfstring{$|\lambda_i|=1$}{|lambda|=1}.}
Anche qui siamo nel caso critico.

\begin{itemize}
    \item Se il blocco di Jordan massimo ha dimensione $1$, i modi restano limitati:
    \[
    1,\qquad \cos(k\theta),\qquad \sin(k\theta).
    \]
    \item Se il blocco ha dimensione $>1$, compaiono modi con fattori $k$, $k^2$, \dots, quindi il sistema diventa polinomialmente divergente.
\end{itemize}

\paragraph{Riassunto per il discreto.}
\begin{itemize}
    \item $|\lambda_i|<1$ $\Rightarrow$ modi convergenti;
    \item $|\lambda_i|>1$ $\Rightarrow$ modi divergenti;
    \item $|\lambda_i|=1$ $\Rightarrow$ bisogna guardare la dimensione del blocco di Jordan massimo.
\end{itemize}

\section{Schema riassuntivo finale}

\subsection*{Tempo continuo}

\begin{center}
\renewcommand{\arraystretch}{1.4}
\begin{tabular}{>{\raggedright\arraybackslash}p{4.2cm} >{\raggedright\arraybackslash}p{9.5cm}}
\toprule
Condizione sugli autovalori & Comportamento dei modi \\
\midrule
$\Re(\lambda_i)<0$ &
Convergenza esponenziale a zero, indipendentemente dalla dimensione dei blocchi di Jordan associati. \\
\addlinespace
$\Re(\lambda_i)>0$ &
Divergenza esponenziale, indipendentemente dalla dimensione dei blocchi di Jordan associati. \\
\addlinespace
$\Re(\lambda_i)=0$ e blocco massimo di dimensione $1$ &
Modi costanti oppure oscillatori limitati e non smorzati. \\
\addlinespace
$\Re(\lambda_i)=0$ e blocco massimo di dimensione $>1$ &
Modi polinomialmente divergenti: compaiono termini del tipo $t, t^2, \dots$. \\
\bottomrule
\end{tabular}
\end{center}

\subsection*{Tempo discreto}

\begin{center}
\renewcommand{\arraystretch}{1.4}
\begin{tabular}{>{\raggedright\arraybackslash}p{4.2cm} >{\raggedright\arraybackslash}p{9.5cm}}
\toprule
Condizione sugli autovalori & Comportamento dei modi \\
\midrule
$|\lambda_i|<1$ &
Convergenza a zero; se $\lambda_i$ è complesso possono esserci oscillazioni smorzate. \\
\addlinespace
$|\lambda_i|>1$ &
Divergenza esponenziale; se $\lambda_i$ è complesso possono esserci oscillazioni amplificate. \\
\addlinespace
$|\lambda_i|=1$ e blocco massimo di dimensione $1$ &
Modi costanti oppure oscillatori a ampiezza costante. \\
\addlinespace
$|\lambda_i|=1$ e blocco massimo di dimensione $>1$ &
Modi polinomialmente divergenti: compaiono termini del tipo $k, k^2, \dots$. \\
\bottomrule
\end{tabular}
\end{center}

\section{Conclusione della lezione}

Il contenuto concettuale della lezione può essere riassunto in una frase:

\begin{quote}
gli andamenti temporali di un sistema lineare sono completamente determinati dagli autovalori e dalla struttura dei blocchi di Jordan associati.
\end{quote}

Più precisamente:
\begin{itemize}
    \item gli autovalori fissano il tipo base di andamento (esponenziale, oscillatorio, crescente, decrescente);
    \item la dimensione dei blocchi di Jordan introduce eventualmente fattori polinomiali in $t$ oppure in $k$;
    \item nei casi critici, questi fattori polinomiali sono decisivi per distinguere tra comportamento limitato e comportamento divergente.
\end{itemize}

Per questo motivo la forma di Jordan non è soltanto uno strumento algebrico: è lo strumento che rende visibile, in modo immediato, la struttura modale del sistema.

\vspace{0.5cm}

Nel seguito del corso, questa idea tornerà continuamente:
\[
\text{spettro + Jordan}
\quad \Longrightarrow \quad
\text{modi}
\quad \Longrightarrow \quad
\text{stabilità, risposta libera, comportamento asintotico}.
\]